\documentclass[9pt,a4paper]{extarticle}
\usepackage[a4paper,margin=9mm,top=8mm,bottom=10mm]{geometry}
% ===== Shared preamble: Class 12 Chemistry (project style, from physics-formula-sheet-source.tex) =====
\usepackage[T1]{fontenc}
\usepackage{lmodern}
\usepackage{amsmath,amssymb}
\usepackage{multicol}
\usepackage[most]{tcolorbox}
\usepackage{enumitem}
\usepackage{booktabs,tabularx,array}
\usepackage{fancyhdr}
\usepackage{xcolor}
\usepackage{microtype}
\usepackage[version=4]{mhchem}
\usepackage{tikz}
\usetikzlibrary{arrows.meta,calc}
\usepackage{pifont}
\usepackage{needspace}
\DeclareUnicodeCharacter{21CC}{\ensuremath{\rightleftharpoons}}
\DeclareUnicodeCharacter{2713}{\ding{51}}
\DeclareUnicodeCharacter{03B1}{\ensuremath{\alpha}}
\DeclareUnicodeCharacter{03B2}{\ensuremath{\beta}}
\DeclareUnicodeCharacter{2192}{\ensuremath{\to}}
\DeclareUnicodeCharacter{2610}{\ensuremath{\square}}
\DeclareUnicodeCharacter{21D2}{\ensuremath{\Rightarrow}}
\DeclareUnicodeCharacter{2194}{\ensuremath{\leftrightarrow}}
\DeclareUnicodeCharacter{2605}{\ding{72}}

% ---------- Colours: one per unit ----------
\definecolor{navy}{HTML}{1B2A4A}
\definecolor{uI}{HTML}{1F5FAD}    % 1 Solutions
\definecolor{uII}{HTML}{0E8A7E}   % 2 Electrochemistry
\definecolor{uIII}{HTML}{C2571A}  % 3 Chemical Kinetics
\definecolor{uIV}{HTML}{6A3FA0}   % 4 d- and f-Block
\definecolor{uV}{HTML}{B0204F}    % 5 Coordination Compounds
\definecolor{uVI}{HTML}{2E8B3A}   % 6 Haloalkanes & Haloarenes
\definecolor{uVII}{HTML}{8A6D00}  % 7 Alcohols, Phenols, Ethers
\definecolor{uVIII}{HTML}{2B6F8F} % 8 Aldehydes, Ketones, Acids
\definecolor{uIX}{HTML}{9A4A1E}   % 9 Amines
\definecolor{uX}{HTML}{5A5A5A}    % 10 Biomolecules
\definecolor{tipY}{HTML}{FFF4B8}
\definecolor{tipYframe}{HTML}{E0B400}
\definecolor{sqp}{HTML}{1B2A4A}
\colorlet{uc}{uI}

% ---------- Tag for questions in the official 2026-27 sample paper ----------
\newcommand{\tagbox}[2]{{\setlength{\fboxsep}{0.8pt}\colorbox{#1}{\textcolor{white}{\sffamily\bfseries\tiny\,#2\,}}}}
\newcommand{\SQP}{\tagbox{sqp}{SQP 26--27}}

% ---------- Unit banner & chapter bar ----------
\newcommand{\unitbanner}[3]{% colour, title, marks
  \colorlet{uc}{#1}%
  \par\vspace{2pt}\noindent
  \begin{tcolorbox}[enhanced,colback=#1,colframe=#1,arc=1.5pt,boxsep=1pt,left=4pt,right=4pt,top=1.5pt,bottom=1.5pt,before skip=2pt,after skip=2pt]
  \sffamily\bfseries\color{white}\footnotesize #2\hfill\scriptsize #3
  \end{tcolorbox}}
\newcommand{\chapbar}[1]{%
  \par\needspace{7\baselineskip}\noindent
  \begin{tcolorbox}[enhanced,colback=uc!18,colframe=uc,boxrule=0pt,leftrule=3pt,arc=0pt,boxsep=0.5pt,left=3pt,right=3pt,top=1pt,bottom=1pt,before skip=3pt,after skip=1.5pt]
  \sffamily\bfseries\color{uc!80!black}\footnotesize #1
  \end{tcolorbox}}
\newcommand{\sub}[1]{\par\needspace{4\baselineskip}\addvspace{2pt}\noindent{\sffamily\bfseries\scriptsize\color{uc}\strut\MakeUppercase{#1}}\par\vspace{0.5pt}}

% ---------- Boxes ----------
\newtcolorbox{kbox}{enhanced,colback=uc!5,colframe=uc,boxrule=0.6pt,arc=1.5pt,left=2pt,right=2pt,top=1pt,bottom=1pt,boxsep=1pt,before skip=2pt,after skip=2pt}
\newtcolorbox{tip}{enhanced,colback=tipY,colframe=tipYframe,boxrule=0.5pt,arc=1.5pt,left=2pt,right=2pt,top=0.5pt,bottom=0.5pt,boxsep=1pt,before skip=2pt,after skip=2pt,
  fontupper=\footnotesize}
\newtcolorbox{numbox}{enhanced,colback=white,colframe=uc!70,boxrule=0.5pt,arc=1.5pt,left=2pt,right=2pt,top=0.5pt,bottom=0.5pt,boxsep=1pt,before skip=2pt,after skip=2pt,
  borderline={0.4pt}{1.2pt}{uc!35,dashed}}
\newcommand{\Tip}{{\sffamily\bfseries\scriptsize TIP\;}}
\newcommand{\Num}{{\sffamily\bfseries\scriptsize\color{uc}STANDARD NUMERICALS\;}}

% ---------- Lists ----------
\setlist{nosep,leftmargin=8pt,labelsep=2pt}
\setlist[itemize]{label={\textcolor{uc}{\tiny$\blacktriangleright$}}}
\newenvironment{fl}{\begin{itemize}}{\end{itemize}}

% ---------- Mini graph helper (tikz) ----------
% \minigraph{width}{height}{x label}{y label}{drawing commands}
\newcommand{\minigraph}[5]{%
\begin{tikzpicture}[x=1cm,y=1cm,line cap=round,font=\tiny]
\draw[-{Stealth[length=3pt]},black!70] (0,0)--(#1,0) node[below left,inner sep=1pt]{#3};
\draw[-{Stealth[length=3pt]},black!70] (0,0)--(0,#2) node[above right,inner sep=1pt]{#4};
#5
\end{tikzpicture}}

\setlength{\columnsep}{9pt}
\setlength{\columnseprule}{0.3pt}
\def\columnseprulecolor{\color{black!25}}
\setlength{\parindent}{0pt}
\setlength{\abovedisplayskip}{1pt}\setlength{\belowdisplayskip}{1pt}
\setlength{\abovedisplayshortskip}{0pt}\setlength{\belowdisplayshortskip}{0pt}
\renewcommand{\arraystretch}{1.08}
\newcolumntype{Y}{>{\raggedright\arraybackslash}X}
\pagestyle{fancy}\fancyhf{}
\renewcommand{\headrulewidth}{0pt}

\fancyfoot[L]{\tiny\sffamily\color{black!55} Class 12 Chemistry · Formula \& Facts Sheet · Part 3: Organic Chemistry · CBSE Board 2026--27 (paper 27 Feb 2027)}
\fancyfoot[R]{\tiny\sffamily\color{black!55} Page \thepage}
\newcommand{\nr}[1]{\textbf{\color{uc!85!black}#1}}
\raggedbottom
\begin{document}
\footnotesize

% ================= TITLE =================
\begin{tcolorbox}[enhanced,colback=navy,colframe=navy,arc=2pt,left=6pt,right=6pt,top=3pt,bottom=3pt]
\color{white}\sffamily
{\Large\bfseries Class 12 Chemistry — Formulas, Reactions \& Key Facts}\hfill{\small NCERT · Rationalised}\\[1pt]
{\footnotesize Part 3 of 3: \textbf{Organic Chemistry} (Ch 6--10) \;·\; CBSE Board 2026--27 only \;·\; Chemistry paper: \textbf{Sat 27 Feb 2027}}\\[2pt]
{\scriptsize Part 1 = Physical (Ch 1--3); Part 2 = Inorganic (Ch 4--5). \SQP\ = asked in the official CBSE 2026--27 sample paper. Ar = aryl (e.g.\ C$_6$H$_5$), R = alkyl.}
\end{tcolorbox}
\vspace{-1pt}
\noindent
\begin{tcolorbox}[enhanced,colback=navy!4,colframe=navy!60,boxrule=0.5pt,arc=1.5pt,left=3pt,right=3pt,top=1pt,bottom=1pt]
\scriptsize\sffamily\textbf{UNIT-WISE MARKS (70)}\rmfamily\quad
1 Solutions 7 · 2 Electrochem 9 · 3 Kinetics 7 · 4 d- \& f-Block 7 · 5 Coordination 7 ·
\textbf{\textcolor{uVI}{6 Haloalkanes 6}} · \textbf{\textcolor{uVII}{7 Alcohols 6}} · \textbf{\textcolor{uVIII}{8 Carbonyls \& Acids 8}} · \textbf{\textcolor{uIX}{9 Amines 6}} · \textbf{\textcolor{uX}{10 Biomolecules 7}}
\;→\; \textbf{Organic = 33 marks (47\%).}\quad
\sffamily\textbf{PAPER}\rmfamily\; A 16×1 (12 MCQ + 4 A--R) · B 5×2 · C 7×3 · D 2×4 case · E 3×5 · 3 h. Organic supplies the most marks: named reactions, conversions, distinguishing tests, ``give reasons'', word problems (identify A, B, C) and a 5-marker.\\
\sffamily\textbf{IN 2026--27 SYLLABUS}\rmfamily\; Ch 6 incl.\ polyhalogen compounds · Ch 7 incl.\ commercially important alcohols · Ch 8 incl.\ uses · Ch 9 incl.\ diazonium salts · Ch 10: carbohydrates, proteins, enzymes, vitamins, nucleic acids, hormones. \textbf{Practical link:} tests for alcoholic, phenolic, aldehydic, ketonic, carboxylic and amino groups; carbohydrates and proteins in foodstuffs.
\end{tcolorbox}

\begin{multicols}{2}\raggedcolumns
% =====================================================================
\unitbanner{uVI}{UNIT 6 · HALOALKANES AND HALOARENES}{6 marks}
\chapbar{Classification, C--X bond, physical properties}
\begin{fl}
\item sp$^3$ C--X: alkyl (1$^\circ$, 2$^\circ$, 3$^\circ$), \textbf{allylic} (X on C next to C=C), \textbf{benzylic} (X on C next to ring). sp$^2$ C--X: \textbf{vinylic} (X on C=C), \textbf{aryl} (X on ring). Geminal (same C) vs vicinal (adjacent C) dihalides.
\item C--X polar (C$^{\delta+}$). Bond length C--F $<$ C--Cl $<$ C--Br $<$ C--I; bond enthalpy reverse. Dipole moment: CH$_3$Cl $>$ CH$_3$F $>$ CH$_3$Br $>$ CH$_3$I.
\item b.p.: RI $>$ RBr $>$ RCl $>$ RF; rises with size; isomers: more branching → lower b.p. (1-bromobutane $>$ 2-bromobutane $>$ 2-bromo-2-methylpropane). p-Dichlorobenzene has highest m.p. of the three isomers (symmetrical, packs well).
\item Only slightly soluble in water: energy to break H-bonds in water is not repaid by new C--X/water attractions. Density rises with mass and number of halogen atoms.
\end{fl}
\chapbar{Preparation}
\begin{kbox}
\ce{R-OH + HCl ->[ZnCl2] R-Cl + H2O} (3$^\circ$ needs no ZnCl$_2$)\quad \ce{3R-OH + PCl3 -> 3R-Cl + H3PO3}\\
\ce{R-OH + PCl5 -> R-Cl + POCl3 + HCl}\quad \ce{R-OH + SOCl2 -> R-Cl + SO2 ^ + HCl ^} (best: by-products escape)\\
Alkene + HX: Markovnikov (H to the C with more H): propene + HBr → 2-bromopropane\\
With peroxide (\textbf{HBr only}): anti-Markovnikov → 1-bromopropane\\
Arenes: \ce{C6H6 + Cl2 ->[Fe / FeCl3][dark] C6H5Cl + HCl} (o/p for substituted rings)
\end{kbox}
\begin{fl}
\item Halogen exchange: \nr{Finkelstein} (NaI, dry acetone), \nr{Swarts} (AgF, Hg$_2$F$_2$, CoF$_2$, SbF$_3$). Aryl halides from diazonium salts: \nr{Sandmeyer}, \nr{Gattermann}, KI. (Equations in the Named-reaction table.)
\item Direct iodination is reversible — needs an oxidant (HNO$_3$, HIO$_4$) to remove HI. Fluorination is too violent.
\end{fl}
\chapbar{Nucleophilic substitution}
{\scriptsize\setlength{\tabcolsep}{2pt}
\begin{tabularx}{\linewidth}{@{}l Y l Y@{}}
\toprule
\textbf{Reagent} & \textbf{Product} & \textbf{Reagent} & \textbf{Product}\\
\midrule
aq.\ KOH/NaOH & alcohol R--OH & KCN (ionic) & nitrile R--C$\equiv$N\\
NaOR$'$ & ether R--O--R$'$ & AgCN (covalent) & isocyanide R--N$\equiv$C\\
NH$_3$ & amines (mixture) & KNO$_2$ & alkyl nitrite R--O--N=O\\
LiAlH$_4$ & alkane R--H & AgNO$_2$ & nitroalkane R--NO$_2$\\
R$'$COOAg & ester R$'$COOR & NaI & R--I\\
\bottomrule
\end{tabularx}\par}
\vspace{1pt}
{\scriptsize\setlength{\tabcolsep}{2pt}
\begin{tabularx}{\linewidth}{@{}l Y Y@{}}
\toprule
 & \textbf{S$_\text{N}$2} & \textbf{S$_\text{N}$1}\\
\midrule
Steps & one step; backside attack via 5-bonded transition state & two: slow ionisation to carbocation, fast attack\\
Rate law & $k$[RX][Nu$^-$] (2nd order) & $k$[RX] (1st order)\\
Reactivity & CH$_3$X $>$ 1$^\circ$ $>$ 2$^\circ$ $>$ 3$^\circ$ (steric hindrance) & 3$^\circ$ $>$ 2$^\circ$ $>$ 1$^\circ$ $>$ CH$_3$X (carbocation stability); allylic, benzylic very fast (resonance)\\
Stereochemistry & \textbf{inversion} (Walden) & \textbf{racemisation} (planar carbocation; attack from both faces)\\
Favoured by & strong nucleophile, aprotic solvent & polar protic solvent (water, alcohol)\\
\bottomrule
\end{tabularx}\par}
\begin{fl}
\item Leaving-group ability (both): R--I $>$ R--Br $>$ R--Cl $>$ R--F.
\item \textbf{Ambident nucleophiles}: \ce{CN-} (C or N) and \ce{NO2-} (O or N).
\item \textbf{Optical activity}: rotates plane-polarised light — dextrorotatory (+, $d$) or laevorotatory ($-$, $l$). \textbf{Chiral} carbon: four different groups (butan-2-ol, 2-bromobutane). Enantiomers = non-superimposable mirror images. \textbf{Racemic mixture} (±): equal amounts of both, optically inactive (external compensation).
\item \textbf{Elimination} (β-elimination) with alc.\ KOH: \nr{Zaitsev rule} — the more substituted alkene is the major product (2-bromobutane → but-2-ene major). 3$^\circ$ halide + strong/bulky base → elimination; 1$^\circ$ + good nucleophile → substitution.
\item \textbf{Metals}: Grignard reagent \ce{R-X + Mg ->[dry ether] RMgX}; reacts with any proton source (\ce{RMgX + H2O -> RH + Mg(OH)X}), so it must be kept dry. \nr{Wurtz}: \ce{2RX + 2Na ->[dry ether] R-R + 2NaX}.
\end{fl}
\chapbar{Haloarenes}
\begin{fl}
\item \textbf{Less reactive to nucleophilic substitution}: (i) resonance gives C--X partial double-bond character; (ii) sp$^2$ C holds the bond tighter (shorter, stronger); (iii) phenyl cation not resonance-stabilised (no S$_\text{N}$1); (iv) electron-rich ring repels nucleophile; backside attack impossible.
\item Chlorobenzene + NaOH, 623 K, 300 atm → phenol (after H$^+$). \textbf{–NO$_2$ at o/p speeds it up} (stabilises the carbanion by resonance): 4-nitrochlorobenzene 443 K; 2,4-dinitro 368 K; 2,4,6-trinitro warm water. –NO$_2$ at meta: no such effect.
\item \textbf{Electrophilic substitution}: halogen is deactivating ($-I$) but \textbf{o/p-directing} ($+R$); para product major. Cl$_2$/FeCl$_3$ → 1,2- + 1,4-dichlorobenzene; conc.\ HNO$_3$/H$_2$SO$_4$ → 1-chloro-2- + 1-chloro-4-nitrobenzene; conc.\ H$_2$SO$_4$ → chlorobenzenesulphonic acids; CH$_3$Cl/AlCl$_3$ → 2- + 4-chlorotoluene; CH$_3$COCl/AlCl$_3$ → 2- + 4-chloroacetophenone.
\item \nr{Wurtz--Fittig}: \ce{ArX + RX + 2Na -> Ar-R + 2NaX}; \nr{Fittig}: \ce{2ArX + 2Na -> Ar-Ar + 2NaX} (biphenyl).
\end{fl}
\chapbar{Polyhalogen compounds — uses \& effects}
{\scriptsize\setlength{\tabcolsep}{2pt}
\begin{tabularx}{\linewidth}{@{}l Y@{}}
\toprule
CH$_2$Cl$_2$ & paint remover, propellant, metal cleaning; harms central nervous system, burns skin, blurs vision\\
CHCl$_3$ & solvent for fats, alkaloids, iodine; was an anaesthetic. Light + air → \textbf{phosgene}: \ce{2CHCl3 + O2 ->[light] 2COCl2 + 2HCl} → stored in closed dark bottles filled to the brim\\
CHI$_3$ & antiseptic (releases free iodine); replaced because of its smell\\
CCl$_4$ & refrigerants, propellants, solvent; in the stratosphere it depletes the ozone layer\\
Freons (CFCs, CCl$_2$F$_2$ = freon-12) & refrigeration, air conditioning, aerosol propellants; very stable, reach the stratosphere and destroy ozone\\
DDT & first chlorinated insecticide (malaria, typhus); not biodegradable, stored in fatty tissue → banned in many countries\\
\bottomrule
\end{tabularx}\par}
\begin{tip}\Tip Markovnikov's exception (peroxide effect) works with HBr only. KCN → nitrile but AgCN → isocyanide; KNO$_2$ → nitrite but AgNO$_2$ → nitro. Allyl/benzyl halides react fastest by S$_\text{N}$1 even though primary. Aryl and vinyl halides do not undergo S$_\text{N}$1/S$_\text{N}$2 under normal conditions.\end{tip}

% =====================================================================
\unitbanner{uVII}{UNIT 7 · ALCOHOLS, PHENOLS AND ETHERS}{6 marks}
\chapbar{Structure \& preparation of alcohols}
\begin{fl}
\item C--O--H angle in methanol 108.9$^\circ$ (lone-pair repulsion). Phenol C--O shorter (136 pm): partial double bond from resonance; sp$^2$ C. Ether C--O--C 111.7$^\circ$ (bulky R groups).
\end{fl}
\begin{kbox}
\textbf{Acid-catalysed hydration} (Markovnikov): \ce{CH3CH=CH2 + H2O ->[H+] CH3CH(OH)CH3}\\
\textbf{Hydroboration--oxidation} (anti-Markovnikov, high yield): propene $\xrightarrow{(\text{BH}_3)_2}\;\xrightarrow{\text{H}_2\text{O}_2/\text{OH}^-}$ propan-1-ol\\
\textbf{Reduction}: aldehyde → 1$^\circ$, ketone → 2$^\circ$ (H$_2$/Pd, NaBH$_4$, LiAlH$_4$); acid → 1$^\circ$ (LiAlH$_4$); commercially via esters (H$_2$/catalyst)\\
\textbf{Grignard}: HCHO → 1$^\circ$; other aldehydes → 2$^\circ$; ketones → 3$^\circ$ (adduct + H$_2$O)
\end{kbox}
\begin{fl}
\item Hydration mechanism: (1) alkene protonated by H$_3$O$^+$ → carbocation; (2) water attacks carbocation; (3) loss of H$^+$ → alcohol.
\end{fl}
\chapbar{Preparation of phenols}
\begin{fl}
\item Chlorobenzene + NaOH (623 K, 300 atm) → sodium phenoxide → H$^+$ → phenol.
\item Benzenesulphonic acid + fused NaOH → sodium phenoxide → H$^+$ → phenol.
\item Diazonium salt + warm water: \ce{C6H5N2+Cl- + H2O -> C6H5OH + N2 + HCl}.
\item \textbf{Cumene} (isopropylbenzene) + O$_2$ → cumene hydroperoxide → H$^+$/H$_2$O → phenol + \textbf{acetone} (by-product).
\end{fl}
\chapbar{Physical properties \& acidity}
\begin{fl}
\item b.p.: alcohols/phenols $\gg$ hydrocarbons, ethers, haloalkanes of similar mass (intermolecular H-bonding). b.p.\ rises with C atoms, falls with branching. Solubility in water (H-bonds with water) falls as the alkyl group grows.
\item Alcohols: \ce{2R-OH + 2Na -> 2R-ONa + H2}. Acidity 1$^\circ$ $>$ 2$^\circ$ $>$ 3$^\circ$ (+I of alkyl groups). Alcohols are \emph{weaker} acids than water: \ce{RO- + H2O -> ROH + OH-}.
\item \textbf{Phenol more acidic than alcohols}: phenoxide ion is resonance-stabilised (negative charge spread over ring); phenol's O on sp$^2$ C. Phenol reacts with NaOH; neither phenol nor alcohol reacts with NaHCO$_3$.
\item –NO$_2$ (EWG) at o/p raises acidity; –CH$_3$ (EDG) lowers it. pK$_a$: 2,4,6-trinitrophenol 0.5 $<$ p-nitrophenol 7.1 $<$ o-nitrophenol 7.2 $<$ m-nitrophenol 8.3 $<$ phenol 10.0 $<$ cresols $\approx$10.2 $<$ ethanol 15.9.
\end{fl}
\chapbar{Reactions of alcohols}
\begin{kbox}
\textbf{Esterification}: \ce{ROH + R$'$COOH <=>[H+] R$'$COOR + H2O}; with R$'$COCl + pyridine (removes HCl)\\
\textbf{Aspirin}: salicylic acid + (CH$_3$CO)$_2$O → acetylsalicylic acid\\
\textbf{Dehydration}: ethanol $\xrightarrow{\text{conc. H}_2\text{SO}_4,\;443\text{ K}}$ ethene; 2$^\circ$: 85\% H$_3$PO$_4$, 440 K; 3$^\circ$: 20\% H$_3$PO$_4$, 358 K. Ease 3$^\circ$ $>$ 2$^\circ$ $>$ 1$^\circ$\\
\textbf{Oxidation}: 1$^\circ$ → aldehyde → acid (acidified KMnO$_4$/K$_2$Cr$_2$O$_7$); stop at aldehyde with \textbf{PCC} or CrO$_3$ (anhydrous); 2$^\circ$ → ketone (CrO$_3$); 3$^\circ$ resist\\
\textbf{Cu, 573 K}: 1$^\circ$ → aldehyde, 2$^\circ$ → ketone, 3$^\circ$ → alkene (dehydration)
\end{kbox}
\begin{fl}
\item With HX/PCl$_3$/PCl$_5$/SOCl$_2$ → R--X. \textbf{Lucas test} (conc.\ HCl + ZnCl$_2$): 3$^\circ$ turbid at once; 2$^\circ$ within minutes; 1$^\circ$ no turbidity at room temperature.
\item \textbf{Dehydration mechanism}: (1) protonation of --OH (fast); (2) loss of H$_2$O → carbocation (slow, rate-determining); (3) loss of H$^+$ from β-C → alkene.
\end{fl}
\chapbar{Reactions of phenol (--OH activates, o/p)}
\begin{fl}
\item \textbf{Nitration}: dilute HNO$_3$, 298 K → o- + p-nitrophenol; separated by steam distillation (o-isomer steam-volatile: intramolecular H-bond; p-isomer: intermolecular H-bond). Conc.\ HNO$_3$ → 2,4,6-trinitrophenol (picric acid).
\item \textbf{Bromination}: Br$_2$ in CS$_2$/CHCl$_3$, low $T$ → o- + p-bromophenol; \textbf{bromine water} → 2,4,6-tribromophenol (white ppt).
\item \nr{Kolbe's reaction} → salicylic acid; \nr{Reimer--Tiemann} → salicylaldehyde (table). Zn dust → benzene. Chromic acid → benzoquinone.
\end{fl}
\chapbar{Commercially important alcohols}
\begin{fl}
\item \textbf{Methanol} (wood spirit): \ce{CO + 2H2 ->[ZnO-Cr2O3][200-300 atm, 573-673 K] CH3OH}. Very poisonous: small amounts cause blindness, larger amounts death. Solvent for paints, varnishes; making HCHO.
\item \textbf{Ethanol}: fermentation of molasses/sugarcane/fruit sugars by yeast enzymes — invertase (sucrose → glucose + fructose), zymase (→ ethanol + CO$_2$), without air. Commercial alcohol is \emph{denatured} with CuSO$_4$ (colour) and pyridine (smell).
\end{fl}
\chapbar{Ethers}
\begin{fl}
\item \textbf{From alcohols}: conc.\ H$_2$SO$_4$, 413 K (excess alcohol) → symmetrical ether. S$_\text{N}$2 mechanism: protonated alcohol attacked by second alcohol molecule. Good only for 1$^\circ$ alcohols (2$^\circ$/3$^\circ$ give alkenes).
\item \nr{Williamson synthesis}: \ce{R-X + R$'$-ONa -> R-O-R$'$ + NaX} (S$_\text{N}$2). Use a \textbf{1$^\circ$ halide} with the bulky alkoxide: \ce{(CH3)3C-ONa + CH3Br} → tert-butyl methyl ether; but \ce{(CH3)3C-Br + CH3ONa} → 2-methylpropene (elimination). Phenoxide + RX → aryl alkyl ether.
\item b.p.\ of ethers $\approx$ alkanes (no H-bonding between molecules); miscibility with water $\approx$ alcohols (H-bond with water).
\item \textbf{Cleavage by HX} (HI $>$ HBr $>$ HCl): \ce{R-O-R + HX -> RX + ROH}. Mixed 1$^\circ$/2$^\circ$: X goes to the smaller group (S$_\text{N}$2). With a 3$^\circ$ group: X goes to the 3$^\circ$ C (S$_\text{N}$1): \ce{CH3-O-C(CH3)3 + HI -> CH3OH + (CH3)3C-I}. \textbf{Anisole} + HI → phenol + CH$_3$I (O--aryl bond stronger; no phenyl cation).
\item Anisole (--OCH$_3$ activating, o/p): Br$_2$ in ethanoic acid → p-bromoanisole (major); CH$_3$Cl/AlCl$_3$ → 2- + 4-methoxytoluene; CH$_3$COCl/AlCl$_3$ → 2- + 4-methoxyacetophenone; HNO$_3$/H$_2$SO$_4$ → 2- + 4-nitroanisole.
\end{fl}
\begin{tip}\Tip Phenol + bromine water gives the \emph{tribromo} product even without a catalyst; in CS$_2$ at low temperature it stops at mono. Never use a 3$^\circ$ alkyl halide in Williamson synthesis. ``Anisole + HI'' never gives iodobenzene.\end{tip}

% =====================================================================
\unitbanner{uVIII}{UNIT 8 · ALDEHYDES, KETONES \& CARBOXYLIC ACIDS}{8 marks}
\chapbar{Carbonyl group \& preparation}
\begin{fl}
\item C=O: sp$^2$ C, trigonal planar ($\approx$120$^\circ$), polar (C$^{\delta+}$, O$^{\delta-}$) → nucleophiles attack C.
\end{fl}
\begin{kbox}
\textbf{Aldehydes}: 1$^\circ$ alcohol + PCC; Cu/573 K; ozonolysis (O$_3$ then Zn/H$_2$O); ethyne + H$_2$O (HgSO$_4$/dil.\ H$_2$SO$_4$) → ethanal;\\
\nr{Rosenmund} (RCOCl + H$_2$, Pd/BaSO$_4$); \nr{Stephen} (RCN, SnCl$_2$/HCl, H$_3$O$^+$); RCN or ester + DIBAL-H, then H$_2$O;\\
toluene → benzaldehyde: \nr{Etard} (CrO$_2$Cl$_2$); CrO$_3$/(CH$_3$CO)$_2$O → benzylidene diacetate → H$_3$O$^+$; Cl$_2$/h$\nu$ → benzal chloride → H$_2$O, 373 K;\\
benzene → benzaldehyde: \nr{Gatterman--Koch} (CO + HCl, anhyd.\ AlCl$_3$/CuCl)\\[1pt]
\textbf{Ketones}: 2$^\circ$ alcohol + K$_2$Cr$_2$O$_7$/H$^+$; alkynes (other than ethyne) + HgSO$_4$/H$_2$SO$_4$;\\
\ce{2RCOCl + R$'$2Cd -> 2RCOR$'$ + CdCl2}; nitrile + Grignard then H$_3$O$^+$ (\ce{CH3CH2CN + C6H5MgBr} → propiophenone); Friedel--Crafts acylation (\ce{C6H6 + RCOCl ->[anhyd. AlCl3] C6H5COR})
\end{kbox}
\begin{fl}
\item b.p.: between hydrocarbons/ethers and alcohols of similar mass (dipole--dipole, no intermolecular H-bonding). Methanal, ethanal, propanone mix with water (H-bond with water).
\end{fl}
\chapbar{Nucleophilic addition (aldehydes $>$ ketones)}
\begin{fl}
\item \textbf{Mechanism}: Nu$^-$ attacks the carbonyl C (perpendicular to the plane) → sp$^3$ tetrahedral alkoxide intermediate → takes H$^+$ → product.
\item Aldehydes more reactive: one alkyl group (less steric hindrance; less +I reduction of C$^{\delta+}$). HCHO $>$ CH$_3$CHO $>$ CH$_3$COCH$_3$. Benzaldehyde $<$ propanal (resonance with ring lowers C$^{\delta+}$).
\end{fl}
{\scriptsize\setlength{\tabcolsep}{2pt}
\begin{tabularx}{\linewidth}{@{}l Y@{}}
\toprule
HCN (trace base) & cyanohydrin \ce{$>$C(OH)CN}\\
NaHSO$_3$ & crystalline bisulphite adduct (separation, purification; regenerates carbonyl with acid/base)\\
RMgX, then H$_2$O & alcohols\\
ROH, dry HCl & aldehyde → hemiacetal → acetal; ketone + ethylene glycol → cyclic ketal\\
\bottomrule
\end{tabularx}\par}
\vspace{1pt}
\textbf{Ammonia derivatives} (\ce{$>$C=O + H2N-Z -> $>$C=N-Z + H2O}, weakly acidic medium):\\
{\scriptsize\setlength{\tabcolsep}{2pt}
\begin{tabularx}{\linewidth}{@{}l l Y@{}}
\toprule
\textbf{Reagent} & \textbf{Z} & \textbf{Product}\\
\midrule
hydroxylamine & --OH & oxime\\
hydrazine & --NH$_2$ & hydrazone\\
phenylhydrazine & --NHC$_6$H$_5$ & phenylhydrazone\\
2,4-DNP & --NHC$_6$H$_3$(NO$_2$)$_2$ & 2,4-dinitrophenylhydrazone (yellow-orange ppt: test for C=O)\\
semicarbazide & --NHCONH$_2$ & semicarbazone\\
\bottomrule
\end{tabularx}\par}
\chapbar{Reduction, oxidation, haloform}
\begin{fl}
\item Reduction: NaBH$_4$/LiAlH$_4$/H$_2$ → alcohols; \nr{Clemmensen} (Zn--Hg, conc.\ HCl) and \nr{Wolff--Kishner} (NH$_2$NH$_2$, then KOH/ethylene glycol, heat) → \ce{$>$CH2}.
\item \textbf{Tollens'} (ammoniacal AgNO$_3$): \ce{RCHO + 2[Ag(NH3)2]+ + 3OH- -> RCOO- + 2Ag + 2H2O + 4NH3} — silver mirror. All aldehydes (aliphatic and aromatic).
\item \textbf{Fehling's} (A: CuSO$_4$; B: alkaline sodium potassium tartrate): \ce{RCHO + 2Cu^2+ + 5OH- -> RCOO- + Cu2O + 3H2O} — red-brown ppt. Aliphatic aldehydes only (benzaldehyde: negative).
\item Ketones need strong oxidants at high $T$ (C--C bonds break → mixture of acids).
\item \textbf{Haloform / iodoform}: CH$_3$CO-- (methyl ketones) and CH$_3$CH(OH)-- compounds + NaOI (I$_2$/NaOH) → yellow CHI$_3$. Positive: ethanal, ethanol, propanone, propan-2-ol, butan-2-one, acetophenone. Negative: methanal, methanol, propanal, pentan-3-one, benzaldehyde.
\end{fl}
\chapbar{α-Hydrogen reactions \& Cannizzaro}
\begin{fl}
\item α-H is acidic: C=O withdraws electrons and the enolate is resonance-stabilised.
\item \nr{Aldol}: carbonyls \textbf{with} α-H + dil.\ NaOH → β-hydroxy aldehyde/ketone → heat → α,β-unsaturated product. \ce{2CH3CHO} → 3-hydroxybutanal → but-2-enal. \ce{2CH3COCH3 ->[Ba(OH)2]} 4-hydroxy-4-methylpentan-2-one → 4-methylpent-3-en-2-one.
\item \textbf{Cross aldol}: two different carbonyls both with α-H → 4 products. One without α-H (benzaldehyde + acetophenone) → mainly benzalacetophenone.
\item \nr{Cannizzaro}: aldehydes \textbf{without} α-H (HCHO, C$_6$H$_5$CHO) + conc.\ alkali → alcohol + carboxylate salt. \ce{2HCHO + conc. KOH -> CH3OH + HCOOK}; \ce{2C6H5CHO + conc. NaOH -> C6H5CH2OH + C6H5COONa}.
\item Ring: --CHO/--CO-- deactivating, \textbf{meta}: benzaldehyde + HNO$_3$/H$_2$SO$_4$ (273--283 K) → 3-nitrobenzaldehyde.
\item \textbf{Uses}: formalin ($\approx$40\% HCHO) preserves specimens; HCHO for bakelite, urea--formaldehyde resins; ethanal → ethanoic acid, ethyl ethanoate; propanone, butanone as solvents; benzaldehyde, vanillin etc.\ in perfumes and flavours.
\end{fl}
\chapbar{Carboxylic acids}
\begin{kbox}
\textbf{Preparation}: 1$^\circ$ alcohol/aldehyde + KMnO$_4$ or K$_2$Cr$_2$O$_7$/H$^+$; alkylbenzene + alk.\ KMnO$_4$ → benzoic acid (whole side chain → --COOH; 3$^\circ$ alkyl not oxidised);\\
nitrile/amide + H$^+$ or OH$^-$ (hydrolysis); \ce{RMgX + CO2 (dry ice)} then H$_3$O$^+$ → RCOOH (\textbf{+1 C}); acyl halides, anhydrides, esters + water (esters + NaOH = saponification)
\end{kbox}
\begin{fl}
\item b.p.\ higher than alcohols of similar mass: stronger H-bonding; exist as \textbf{dimers} in vapour and aprotic solvents. C$_1$--C$_4$ acids mix with water.
\item \textbf{Acidity}: carboxylate ion has two \emph{equivalent} resonance structures (charge on two O atoms) → RCOOH more acidic than phenol and alcohols. RCOOH + NaHCO$_3$ → CO$_2$ effervescence (phenol does not).
\item EWG increase acidity, EDG decrease. Nearer the EWG to --COOH, stronger the acid. Phenyl/vinyl directly on --COOH increases acidity (sp$^2$ C).
\end{fl}
{\scriptsize\textbf{NCERT acidity order:} CF$_3$COOH $>$ CCl$_3$COOH $>$ CHCl$_2$COOH $>$ NO$_2$CH$_2$COOH $>$ NCCH$_2$COOH $>$ FCH$_2$COOH $>$ ClCH$_2$COOH $>$ BrCH$_2$COOH $>$ HCOOH $>$ ClCH$_2$CH$_2$COOH $>$ C$_6$H$_5$COOH $>$ C$_6$H$_5$CH$_2$COOH $>$ CH$_3$COOH $>$ CH$_3$CH$_2$COOH.\par}
\begin{kbox}
\textbf{Reactions}: 2RCOOH $\xrightarrow{\text{H}^+,\ \text{P}_2\text{O}_5,\ \Delta}$ anhydride;\quad esterification (H$^+$, reversible);\\
\ce{RCOOH + SOCl2 -> RCOCl + SO2 + HCl} (also PCl$_5$, PCl$_3$);\quad \ce{RCOOH + NH3 -> RCOONH4 ->[\Delta] RCONH2} (phthalic acid → phthalimide);\\
LiAlH$_4$ or B$_2$H$_6$ → RCH$_2$OH (NaBH$_4$ does not reduce --COOH; B$_2$H$_6$ leaves ester, nitro, halo groups alone);\\
\textbf{Decarboxylation}: RCOONa + soda lime (NaOH + CaO), heat → R--H (\textbf{$-$1 C});\quad \nr{Kolbe electrolysis}: aqueous RCOONa → R--R;\\
\nr{HVZ}: \ce{RCH2COOH ->[X2 / red P][H2O] RCH(X)COOH} (α-halo acid);\quad ring: --COOH deactivating, meta (benzoic acid → 3-nitro/3-bromobenzoic acid); \textbf{no Friedel--Crafts} (AlCl$_3$ binds to --COOH)
\end{kbox}
\begin{fl}
\item Uses: HCOOH (rubber, textile, dyeing, leather, electroplating); CH$_3$COOH (vinegar, solvent); hexanedioic acid (nylon-6,6); benzoate esters (perfumery); sodium benzoate (food preservative); higher fatty acids (soaps, detergents).
\end{fl}

% =====================================================================
\unitbanner{uIX}{UNIT 9 · AMINES}{6 marks}
\chapbar{Structure, naming, preparation}
\begin{fl}
\item N is sp$^3$, pyramidal (lone pair); C--N--C in trimethylamine 108$^\circ$. 1$^\circ$ RNH$_2$, 2$^\circ$ R$_2$NH, 3$^\circ$ R$_3$N, quaternary R$_4$N$^+$X$^-$.
\item IUPAC: ethanamine; N-methylmethanamine; N,N-dimethylmethanamine; benzenamine (aniline); prop-2-en-1-amine; hexane-1,6-diamine.
\end{fl}
\begin{kbox}
Nitro → amine: H$_2$/Ni, Pd or Pt; Sn/HCl or \textbf{Fe/HCl} (preferred: FeCl$_2$ hydrolyses and regenerates HCl, so little acid needed)\\
Ammonolysis: \ce{RX + NH3 ->[ethanol][373 K, sealed tube]} \ce{RNH3+X-} → 1$^\circ$, 2$^\circ$, 3$^\circ$ amines + quaternary salt (\textbf{mixture}; excess NH$_3$ favours 1$^\circ$). RI $>$ RBr $>$ RCl\\
\ce{RC#N ->[LiAlH4 / H2-Ni / Na(Hg)-C2H5OH] RCH2NH2} (\textbf{+1 C});\quad \ce{RCONH2 ->[LiAlH4][H2O] RCH2NH2}\\
\nr{Gabriel phthalimide}: phthalimide + alc.\ KOH → K-phthalimide + RX → N-alkylphthalimide → NaOH(aq) → RNH$_2$ (1$^\circ$ aliphatic only)\\
\nr{Hoffmann bromamide}: \ce{RCONH2 + Br2 + 4NaOH -> RNH2 + Na2CO3 + 2NaBr + 2H2O} (\textbf{$-$1 C})
\end{kbox}
\begin{fl}
\item Gabriel cannot make aniline: aryl halides do not undergo nucleophilic substitution with the phthalimide anion.
\item b.p.: 1$^\circ$ $>$ 2$^\circ$ $>$ 3$^\circ$ (3$^\circ$ has no N--H for H-bonding); amines $<$ alcohols of similar mass (N less electronegative than O). Lower amines: fishy smell. Aniline darkens on storage (air oxidation).
\end{fl}
\chapbar{Basicity}
\begin{fl}
\item Lone pair on N accepts H$^+$; larger $K_b$ / smaller p$K_b$ = stronger base. Aliphatic amines $>$ NH$_3$ (+I effect). \textbf{Aniline $<$ NH$_3$}: lone pair delocalised into the ring; anilinium ion has fewer resonance structures than aniline.
\item Gas phase: 3$^\circ$ $>$ 2$^\circ$ $>$ 1$^\circ$ $>$ NH$_3$ (+I only).
\item Water (+I, solvation by H-bonding, steric): \textbf{(CH$_3$)$_2$NH $>$ CH$_3$NH$_2$ $>$ (CH$_3$)$_3$N $>$ NH$_3$}; \textbf{(C$_2$H$_5$)$_2$NH $>$ (C$_2$H$_5$)$_3$N $>$ C$_2$H$_5$NH$_2$ $>$ NH$_3$}.
\item p$K_b$: (C$_2$H$_5$)$_2$NH 3.00, (CH$_3$)$_2$NH 3.27, C$_2$H$_5$NH$_2$ 3.29, CH$_3$NH$_2$ 3.38, NH$_3$ 4.75, benzylamine 4.70, aniline 9.38.
\item Substituted anilines: EDG (--CH$_3$, --OCH$_3$) raise basicity; EWG (--NO$_2$, --COOH, --SO$_3$H, --X) lower it. p-toluidine $>$ aniline $>$ p-nitroaniline. Benzylamine $>$ aniline (N not attached to ring).
\end{fl}
\chapbar{Reactions of amines}
\begin{fl}
\item \textbf{Acylation} (RCOCl, anhydride; pyridine removes HCl) → amides. \textbf{Benzoylation} with C$_6$H$_5$COCl: \ce{CH3NH2} → N-methylbenzamide.
\item \nr{Carbylamine}: \ce{R-NH2 + CHCl3 + 3KOH ->[\Delta] R-NC + 3KCl + 3H2O} — foul-smelling isocyanide; 1$^\circ$ amines only (aliphatic and aromatic).
\item \textbf{Nitrous acid} (NaNO$_2$ + HCl): 1$^\circ$ aliphatic → alcohol + \textbf{N$_2$} (quantitative; estimation of amino acids/proteins). Aniline at 273--278 K → diazonium salt.
\item \nr{Hinsberg test} (benzenesulphonyl chloride): 1$^\circ$ → sulphonamide \textbf{soluble in alkali} (acidic N--H); 2$^\circ$ → sulphonamide \textbf{insoluble in alkali} (no N--H); 3$^\circ$ → no reaction.
\item \textbf{Aniline + bromine water} → 2,4,6-tribromoaniline (white ppt). For mono: acetylate (acetanilide, less activating: N lone pair also delocalised onto C=O) → brominate → 4-bromoacetanilide → hydrolyse → 4-bromoaniline.
\item \textbf{Nitration}: direct gives tarry oxidation products and 47\% meta (anilinium ion is meta-directing); 51\% p, 2\% o. Protect by acetylation → p-nitroaniline.
\item \textbf{Sulphonation}: conc.\ H$_2$SO$_4$ → anilinium hydrogensulphate → 453--473 K → sulphanilic acid (zwitterion).
\item \textbf{No Friedel--Crafts}: aniline forms a salt with AlCl$_3$; N$^+$ strongly deactivates the ring.
\end{fl}
\chapbar{Diazonium salts}
\begin{kbox}
Diazotisation: \ce{C6H5NH2 + NaNO2 + 2HCl ->[273-278 K] C6H5N2+Cl- + NaCl + 2H2O}\\
Aryl diazonium salts are stable at low $T$ (resonance); used at once. Alkyl diazonium salts decompose immediately. \ce{C6H5N2+BF4-} is stable at room temperature.
\end{kbox}
{\scriptsize\setlength{\tabcolsep}{2pt}
\begin{tabularx}{\linewidth}{@{}l Y@{}}
\toprule
\textbf{Reagent (N$_2$ lost)} & \textbf{Product}\\
\midrule
CuCl/HCl; CuBr/HBr (\nr{Sandmeyer}) & ArCl; ArBr\\
Cu powder + HCl/HBr (\nr{Gattermann}) & ArCl; ArBr\\
CuCN/KCN & ArCN (benzonitrile)\\
KI & ArI\\
HBF$_4$, then heat & ArF (+ BF$_3$ + N$_2$)\\
H$_3$PO$_2$ or CH$_3$CH$_2$OH & ArH (benzene)\\
H$_2$O, warm & ArOH (phenol)\\
HBF$_4$, then NaNO$_2$/Cu, heat & ArNO$_2$\\
\midrule
\textbf{Coupling (N retained)} & \\
phenol, alkaline (pH 9--10) & p-hydroxyazobenzene (\textbf{orange} dye)\\
aniline, mild acid & p-aminoazobenzene (\textbf{yellow} dye)\\
\bottomrule
\end{tabularx}\par}
\begin{fl}
\item \textbf{Importance}: F, I, CN, OH can be put on the ring this way, which direct halogenation/substitution cannot do (ArF, ArI, ArCN, 1,3,5-tribromobenzene via aniline).
\end{fl}
\begin{tip}\Tip Aniline is \emph{o/p}-directing, but anilinium (in acid) is \emph{meta}-directing. Carbylamine = test for 1$^\circ$ amines only. Hoffmann bromamide removes a C; nitrile reduction adds a C. Diazonium salts: 0--5 $^\circ$C (273--278 K).\end{tip}

% =====================================================================
\unitbanner{uX}{UNIT 10 · BIOMOLECULES}{7 marks}
\chapbar{Carbohydrates}
\begin{fl}
\item Optically active polyhydroxy aldehydes or ketones, or compounds that give these on hydrolysis. Monosaccharides (glucose, fructose, ribose); oligosaccharides (2--10 units: sucrose, maltose, lactose); polysaccharides (starch, cellulose, glycogen). Aldoses vs ketoses; glucose = aldohexose, fructose = ketohexose.
\item \textbf{Reducing sugars} (reduce Tollens'/Fehling's — free --CHO/C=O): all monosaccharides, maltose, lactose. \textbf{Non-reducing}: sucrose.
\item Glucose from sucrose (dil.\ HCl/H$_2$SO$_4$, warm) or starch (dil.\ H$_2$SO$_4$, 393 K, 2--3 atm).
\end{fl}
{\scriptsize\setlength{\tabcolsep}{2pt}
\begin{tabularx}{\linewidth}{@{}l Y@{}}
\toprule
\textbf{Glucose + reagent} & \textbf{Result → conclusion}\\
\midrule
HI, prolonged heating & n-hexane → 6 C in a straight chain\\
NH$_2$OH; HCN & oxime; cyanohydrin → carbonyl group present\\
Bromine water & gluconic acid (C$_6$) → carbonyl is an \textbf{aldehyde}\\
Acetic anhydride & glucose pentaacetate → \textbf{five --OH} groups on different C\\
Nitric acid & saccharic acid (dicarboxylic) → a \textbf{1$^\circ$ alcohol} group present\\
\bottomrule
\end{tabularx}\par}
\begin{fl}
\item \textbf{Open chain fails to explain}: no 2,4-DNP or Schiff's test and no NaHSO$_3$ adduct; pentaacetate does not react with NH$_2$OH; two forms α (m.p.\ 419 K) and β (423 K).
\item \textbf{Cyclic}: C5--OH adds to --CHO → six-membered \textbf{pyranose} ring (hemiacetal). α and β differ only at C1 — \textbf{anomers}; C1 = anomeric carbon. Haworth projection.
\item D/L: configuration relative to glyceraldehyde (--OH on right at the lowest chiral C = D); (+)/($-$) = direction of rotation. Glucose: D-(+); fructose: D-($-$), five-membered \textbf{furanose} ring (C5--OH + C2 carbonyl).
\item \textbf{Glycosidic linkage}: oxide link joining two monosaccharide units.
\end{fl}
{\scriptsize\setlength{\tabcolsep}{2pt}
\begin{tabularx}{\linewidth}{@{}l Y l@{}}
\toprule
\textbf{Sugar} & \textbf{Units \& link} & \textbf{Reducing?}\\
\midrule
Sucrose & α-D-glucose C1 -- β-D-fructose C2 & no\\
Maltose & α-D-glucose C1 -- α-D-glucose C4 & yes\\
Lactose (milk) & β-D-galactose C1 -- β-D-glucose C4 & yes\\
Starch & α-glucose: amylose (15--20\%, soluble, unbranched, C1--C4); amylopectin (80--85\%, insoluble, branched, C1--C4 + C1--C6) & --\\
Cellulose & β-D-glucose, straight chains, C1--C4 (plant cell walls; humans cannot digest) & --\\
Glycogen & like amylopectin, more branched; animal starch (liver, muscles, brain) & --\\
\bottomrule
\end{tabularx}\par}
\begin{fl}
\item \textbf{Invert sugar}: sucrose (dextrorotatory) → glucose ($+52.5^\circ$) + fructose ($-92.4^\circ$): mixture is laevorotatory.
\item Importance: energy source; storage (starch, glycogen); cell walls (cellulose); wood, cotton; ribose/deoxyribose in nucleic acids.
\end{fl}
\chapbar{Proteins}
\begin{fl}
\item α-Amino acids \ce{H2N-CH(R)-COOH}. \textbf{Essential} (from diet): valine, leucine, isoleucine, arginine, lysine, threonine, methionine, phenylalanine, tryptophan, histidine. Neutral/acidic/basic by number of --NH$_2$ and --COOH.
\item \textbf{Zwitterion} \ce{H3N+-CH(R)-COO-}: amphoteric, high m.p., dipolar salt. All except glycine are optically active, mostly L.
\item \textbf{Peptide bond} --CO--NH--: amide link between --COOH of one and --NH$_2$ of the next, with loss of water. $>$100 residues, mass $>$10\,000 u → protein.
\item \textbf{Fibrous} (parallel chains, H-bonds/disulphide links; insoluble): keratin (hair, wool, silk), myosin (muscle). \textbf{Globular} (coiled spheres; soluble): insulin, albumins.
\item Structure: \textbf{primary} — sequence; \textbf{secondary} — α-helix (right-handed, H-bonds between --NH and C=O of adjacent turns) or β-pleated sheet (chains side by side, H-bonds); \textbf{tertiary} — further folding (H-bonds, disulphide, van der Waals, electrostatic); \textbf{quaternary} — arrangement of two or more subunits.
\item \textbf{Denaturation}: heat or pH change breaks H-bonds; protein unfolds and loses activity; 2$^\circ$/3$^\circ$ lost, \textbf{primary intact}. Egg white coagulates on boiling; milk curdles (lactic acid).
\item \textbf{Enzymes}: biocatalysts, mostly globular proteins, highly specific (maltase → maltose; sucrase/invertase → sucrose); lower activation energy.
\end{fl}
\chapbar{Vitamins}
{\scriptsize\setlength{\tabcolsep}{2pt}
\begin{tabularx}{\linewidth}{@{}l Y Y@{}}
\toprule
\textbf{Vitamin} & \textbf{Source} & \textbf{Deficiency}\\
\midrule
A & fish liver oil, carrots, butter, milk & xerophthalmia (hardening of cornea), night blindness\\
B$_1$ (thiamine) & yeast, milk, green vegetables, cereals & beriberi\\
B$_2$ (riboflavin) & milk, egg white, liver, kidney & cheilosis, digestive disorders, burning skin\\
B$_6$ (pyridoxine) & yeast, milk, egg yolk, cereals, grams & convulsions\\
B$_{12}$ & meat, fish, egg, curd & pernicious anaemia\\
C (ascorbic acid) & citrus fruits, amla, green leafy vegetables & scurvy (bleeding gums)\\
D & sunlight, fish, egg yolk & rickets (children), osteomalacia (adults)\\
E & vegetable oils (wheat germ, sunflower) & fragile RBCs, muscular weakness\\
K & green leafy vegetables & longer blood-clotting time\\
\bottomrule
\end{tabularx}\par}
\begin{fl}
\item Fat-soluble A, D, E, K (stored in liver, adipose tissue). Water-soluble B group and C — excreted in urine, must be taken regularly (except B$_{12}$).
\end{fl}
\chapbar{Nucleic acids \& hormones}
\begin{fl}
\item Nucleotide = base + pentose sugar + phosphate; nucleoside = base + sugar (base at C1$'$; phosphate at C5$'$). Nucleotides joined by \textbf{phosphodiester} links (C5$'$--C3$'$).
\item \textbf{DNA}: β-D-2-deoxyribose; bases A, G (purines), C, T; double helix (Watson--Crick), A=T (2 H-bonds), G$\equiv$C (3 H-bonds); heredity, replication. \textbf{RNA}: β-D-ribose; A, G, C, \textbf{U}; usually single strand; mRNA, rRNA, tRNA; protein synthesis.
\item \textbf{Hormones}: chemical messengers from endocrine glands. Steroids (oestrogens, androgens), polypeptides (insulin, endorphins), amino-acid derivatives (adrenaline, noradrenaline, thyroxine). Insulin + glucagon keep blood glucose steady; low thyroxine → hypothyroidism (lethargy, obesity), excess → hyperthyroidism; low iodine → goitre. Glucocorticoids (carbohydrate metabolism, inflammation), mineralocorticoids (water/salt excretion); testosterone, oestradiol, progesterone.
\end{fl}
\begin{tip}\Tip Sucrose is non-reducing because both anomeric carbons are tied up in the glycosidic bond. Cellulose is β-linked (indigestible), starch α-linked. DNA has thymine, RNA uracil. Denaturation keeps the primary structure. Vitamin C cannot be stored.\end{tip}
\end{multicols}

% =====================================================================
\clearpage
\begin{tcolorbox}[enhanced,colback=navy,colframe=navy,arc=2pt,left=6pt,right=6pt,top=2pt,bottom=2pt,before skip=4pt]
\color{white}\sffamily\bfseries\small NAMED REACTIONS — reagents, conditions, product \hfill\footnotesize\mdseries the most-asked 2-/3-mark item in organic
\end{tcolorbox}
{\scriptsize\setlength{\tabcolsep}{2.5pt}\renewcommand{\arraystretch}{1.18}
\begin{tabularx}{\textwidth}{@{}>{\raggedright\arraybackslash}p{2.7cm} Y p{0.1cm}>{\raggedright\arraybackslash}p{2.7cm} Y @{}}
\toprule
\textbf{Name} & \textbf{Equation / conditions} & & \textbf{Name} & \textbf{Equation / conditions}\\
\midrule
\textcolor{uVI}{\textbf{Finkelstein}} & \ce{R-X + NaI ->[dry acetone] R-I + NaX} (X = Cl, Br; NaX precipitates) & &
\textcolor{uVIII}{\textbf{Etard}} & toluene \ce{->[CrO2Cl2, CS2]} chromium complex \ce{->[H3O+]} benzaldehyde\\
\textcolor{uVI}{\textbf{Swarts}} & \ce{R-Br + AgF -> R-F + AgBr} (also Hg$_2$F$_2$, CoF$_2$, SbF$_3$) & &
\textcolor{uVIII}{\textbf{Gatterman--Koch}} & \ce{C6H6 + CO + HCl ->[anhyd. AlCl3 / CuCl] C6H5CHO}\\
\textcolor{uVI}{\textbf{Sandmeyer}} & \ce{ArN2+X- ->[CuCl/HCl or CuBr/HBr] ArCl/ArBr + N2}; CuCN/KCN → ArCN & &
\textcolor{uVIII}{\textbf{Clemmensen}} & \ce{$>$C=O ->[Zn-Hg, conc. HCl] $>$CH2 + H2O}\\
\textcolor{uVI}{\textbf{Gattermann}} & \ce{ArN2+X- ->[Cu / HX] ArX + N2} & &
\textcolor{uVIII}{\textbf{Wolff--Kishner}} & \ce{$>$C=O ->[NH2NH2] $>$C=NNH2 ->[KOH / ethylene glycol][\Delta] $>$CH2 + N2}\\
\textcolor{uVI}{\textbf{Wurtz}} & \ce{2R-X + 2Na ->[dry ether] R-R + 2NaX} & &
\textcolor{uVIII}{\textbf{Aldol}} & \ce{2CH3CHO ->[dil. NaOH] CH3CH(OH)CH2CHO ->[\Delta] CH3CH=CHCHO + H2O}\\
\textcolor{uVI}{\textbf{Wurtz--Fittig}} & \ce{ArX + RX + 2Na ->[dry ether] Ar-R + 2NaX} & &
\textcolor{uVIII}{\textbf{Cannizzaro}} & \ce{2HCHO ->[conc. KOH] CH3OH + HCOOK} (no α-H)\\
\textcolor{uVI}{\textbf{Fittig}} & \ce{2ArX + 2Na ->[dry ether] Ar-Ar + 2NaX} & &
\textcolor{uVIII}{\textbf{HVZ}} & \ce{RCH2COOH ->[(i) X2 / red P][(ii) H2O] RCH(X)COOH}\\
\textcolor{uVI}{\textbf{Friedel--Crafts}} & alkylation: \ce{ArH + RCl ->[anhyd. AlCl3] ArR + HCl}; acylation: \ce{ArH + RCOCl ->[anhyd. AlCl3] ArCOR + HCl} & &
\textcolor{uVIII}{\textbf{Kolbe electrolysis}} & \ce{2RCOO^- Na+ + 2H2O ->[electrolysis] R-R + 2CO2 + H2 + 2NaOH}\\
\textcolor{uVII}{\textbf{Kolbe's reaction}} & sodium phenoxide \ce{->[CO2][400 K, 4-7 atm]} \ce{->[H+]} salicylic acid (2-hydroxybenzoic acid) & &
\textcolor{uVIII}{\textbf{Decarboxylation}} & \ce{RCOONa ->[NaOH + CaO][\Delta] R-H + Na2CO3}\\
\textcolor{uVII}{\textbf{Reimer--Tiemann}} & phenol \ce{->[CHCl3 + aq. NaOH][340 K]} (benzal chloride intermediate) \ce{->[H+]} salicylaldehyde & &
\textcolor{uIX}{\textbf{Gabriel}} & phthalimide \ce{->[alc. KOH]} K-salt \ce{->[RX]} N-alkylphthalimide \ce{->[NaOH(aq)]} \ce{RNH2} + sodium phthalate\\
\textcolor{uVII}{\textbf{Williamson}} & \ce{R-X + R$'$-O^- Na+ -> R-O-R$'$ + NaX} (1$^\circ$ R--X) & &
\textcolor{uIX}{\textbf{Hoffmann bromamide}} & \ce{RCONH2 + Br2 + 4NaOH -> RNH2 + Na2CO3 + 2NaBr + 2H2O}\\
\textcolor{uVIII}{\textbf{Rosenmund}} & \ce{RCOCl + H2 ->[Pd / BaSO4] RCHO + HCl} (poisoned catalyst stops at aldehyde) & &
\textcolor{uIX}{\textbf{Carbylamine}} & \ce{RNH2 + CHCl3 + 3KOH ->[\Delta] RNC + 3KCl + 3H2O}\\
\textcolor{uVIII}{\textbf{Stephen}} & \ce{RCN + SnCl2 + HCl -> RCH=NH ->[H3O+] RCHO} & &
\textcolor{uIX}{\textbf{Diazotisation / coupling}} & \ce{ArNH2 ->[NaNO2 + HCl][273-278 K] ArN2+Cl-}; + phenol (OH$^-$) → orange azo dye; + aniline (H$^+$) → yellow\\
\bottomrule
\end{tabularx}\par}

\begin{multicols}{2}\raggedcolumns
\colorlet{uc}{navy}
\unitbanner{navy}{DISTINGUISHING TESTS}{reagent → what the positive one shows}
{\scriptsize\setlength{\tabcolsep}{2pt}
\begin{tabularx}{\linewidth}{@{}>{\raggedright\arraybackslash}p{0.33\linewidth} Y@{}}
\toprule
\textbf{Pair} & \textbf{Test (positive result)}\\
\midrule
1$^\circ$ / 2$^\circ$ / 3$^\circ$ alcohol & Lucas: 3$^\circ$ turbid at once, 2$^\circ$ in minutes, 1$^\circ$ none\\
Phenol vs alcohol & neutral FeCl$_3$: violet (phenol); Br$_2$ water: white ppt (phenol); dissolves in NaOH (phenol)\\
Ethanol vs methanol & iodoform: yellow CHI$_3$ (ethanol)\\
Propan-2-ol vs propan-1-ol & iodoform (propan-2-ol); Lucas (2$^\circ$ faster)\\
Aldehyde vs ketone & Tollens': silver mirror; Fehling's: red ppt (aliphatic aldehydes)\\
Propanal vs propanone & Tollens'/Fehling's (propanal); iodoform (propanone)\\
Benzaldehyde vs ethanal & Fehling's and iodoform (ethanal positive)\\
Benzaldehyde vs acetophenone & Tollens' (benzaldehyde); iodoform (acetophenone)\\
Pentan-2-one vs pentan-3-one & iodoform (pentan-2-one)\\
Carboxylic acid vs phenol & NaHCO$_3$: CO$_2$ effervescence (acid)\\
Benzoic acid vs ethyl benzoate & NaHCO$_3$ (acid)\\
1$^\circ$ / 2$^\circ$ / 3$^\circ$ amine & Hinsberg: 1$^\circ$ product soluble in alkali; 2$^\circ$ insoluble; 3$^\circ$ no reaction\\
1$^\circ$ amine vs 2$^\circ$/3$^\circ$ & carbylamine: foul smell (1$^\circ$)\\
Aniline vs ethylamine / benzylamine & NaNO$_2$/HCl at 0--5 $^\circ$C, then β-naphthol (alkaline): orange-red dye (aniline); aliphatic gives N$_2$ bubbles\\
Aniline vs N-methylaniline & carbylamine (aniline)\\
Glucose vs sucrose & Tollens'/Fehling's (glucose)\\
C=O compound present & 2,4-DNP: yellow-orange ppt\\
\bottomrule
\end{tabularx}\par}

\unitbanner{navy}{ORDERS TO KNOW}{write the reason with the order}
{\scriptsize
\begin{fl}
\item S$_\text{N}$2: CH$_3$X $>$ 1$^\circ$ $>$ 2$^\circ$ $>$ 3$^\circ$ \;·\; S$_\text{N}$1: 3$^\circ$ $>$ 2$^\circ$ $>$ 1$^\circ$ $>$ CH$_3$X \;·\; leaving group I $>$ Br $>$ Cl $>$ F.
\item b.p.\ (similar mass): carboxylic acid $>$ alcohol $>$ 1$^\circ$ amine $>$ aldehyde/ketone $>$ ether $\approx$ alkane. Amines 1$^\circ$ $>$ 2$^\circ$ $>$ 3$^\circ$. Haloalkanes RI $>$ RBr $>$ RCl.
\item Acidity: RCOOH $>$ phenol $>$ water $>$ alcohol; alcohols 1$^\circ$ $>$ 2$^\circ$ $>$ 3$^\circ$; picric acid $>$ p-nitrophenol $>$ o- $>$ m-nitrophenol $>$ phenol $>$ cresol.
\item Carboxylic acids: F-- $>$ Cl-- $>$ Br--CH$_2$COOH; CCl$_3$ $>$ CHCl$_2$ $>$ CH$_2$ClCOOH; 2-chloro $>$ 3-chloro $>$ 4-chlorobutanoic; HCOOH $>$ CH$_3$COOH $>$ CH$_3$CH$_2$COOH; 4-nitrobenzoic $>$ benzoic $>$ 4-methylbenzoic.
\item Nucleophilic addition: HCHO $>$ CH$_3$CHO $>$ CH$_3$COCH$_3$; aliphatic aldehyde $>$ benzaldehyde.
\item Basicity (water): (CH$_3$)$_2$NH $>$ CH$_3$NH$_2$ $>$ (CH$_3$)$_3$N $>$ NH$_3$ $>$ C$_6$H$_5$NH$_2$; (C$_2$H$_5$)$_2$NH $>$ (C$_2$H$_5$)$_3$N $>$ C$_2$H$_5$NH$_2$; benzylamine $>$ aniline; p-toluidine $>$ aniline $>$ p-nitroaniline.
\item Reactivity of alcohols with HX / dehydration: 3$^\circ$ $>$ 2$^\circ$ $>$ 1$^\circ$; HI $>$ HBr $>$ HCl.
\end{fl}\par}

\unitbanner{navy}{CONVERSIONS — STEP CHAINS}{carbon count first}
{\scriptsize
\textbf{Step up one C}: R--X + KCN → RCN (then LiAlH$_4$ → amine, or H$_3$O$^+$ → acid); RMgX + CO$_2$ → RCOOH; RMgX + HCHO → RCH$_2$OH. \textbf{Step down one C}: Hoffmann bromamide; decarboxylation. \textbf{Double}: Wurtz, Kolbe electrolysis.
\begin{fl}
\item Methanamine → ethanamine: HNO$_2$ → CH$_3$OH → PCl$_5$ → CH$_3$Cl → KCN → CH$_3$CN → LiAlH$_4$.
\item Ethanamine → methanamine: HNO$_2$ → C$_2$H$_5$OH → K$_2$Cr$_2$O$_7$/H$^+$ → CH$_3$COOH → NH$_3$, heat → CH$_3$CONH$_2$ → Br$_2$/NaOH.
\item Ethanoic acid → methanamine: NH$_3$, heat → acetamide → Br$_2$/NaOH.
\item Benzene → aniline: conc.\ HNO$_3$/H$_2$SO$_4$ → Sn/HCl. Aniline → chlorobenzene / benzonitrile / phenol / benzene / fluorobenzene / iodobenzene: diazotise, then CuCl / CuCN / H$_2$O / H$_3$PO$_2$ / HBF$_4$ + heat / KI.
\item Aniline → benzoic acid: diazotise → CuCN → H$_3$O$^+$. Aniline → p-bromoaniline: acetylate → Br$_2$ → hydrolyse.
\item Propene → propan-1-ol: B$_2$H$_6$, H$_2$O$_2$/OH$^-$. Propene → propan-2-ol: H$_2$O/H$^+$. Propanone → propene: NaBH$_4$ → conc.\ H$_2$SO$_4$, heat.
\item Benzoic acid → benzaldehyde: SOCl$_2$ → Rosenmund. Benzene → acetophenone: CH$_3$COCl/AlCl$_3$. Acetophenone → benzoic acid: NaOI (haloform) → H$^+$.
\item Bromomethane → propanone: Mg/ether → CH$_3$MgBr + CH$_3$CN → H$_3$O$^+$.
\item Ethanal → but-2-enal: dil.\ NaOH (aldol), heat. Ethanol → 3-hydroxybutanal: PCC → ethanal → dil.\ NaOH.
\item Phenol → aspirin: Kolbe → salicylic acid → (CH$_3$CO)$_2$O. Phenol → anisole: NaOH → CH$_3$I. Chlorobenzene → phenol: NaOH, 623 K, 300 atm → H$^+$.
\item Ethanol → propanenitrile: PCl$_5$ → KCN. Toluene → benzoic acid: alk.\ KMnO$_4$ → H$^+$.
\end{fl}\par}

\unitbanner{navy}{TOP ``GIVE REASON'' ANSWERS}{one cause, stated in chemistry terms}
{\scriptsize
\begin{fl}
\item \textbf{Haloarenes resist nucleophilic substitution} — resonance (partial C=X), sp$^2$ C, unstable phenyl cation, electron-rich ring repels Nu$^-$.
\item \textbf{Grignard reagents are made in dry ether} — they react with water to give alkanes.
\item \textbf{Chloroform stored in dark bottles filled to the brim} — air + light oxidise it to poisonous phosgene.
\item \textbf{Phenol is acidic, ethanol hardly} — phenoxide ion is resonance-stabilised; ethoxide is not (and is destabilised by +I).
\item \textbf{o-Nitrophenol is steam-volatile, p-isomer not} — intramolecular vs intermolecular H-bonding.
\item \textbf{Alcohols boil higher than ethers of similar mass} — intermolecular H-bonding.
\item \textbf{Anisole + HI gives phenol + CH$_3$I} — O--C(aryl) has partial double-bond character; S$_\text{N}$2 attack occurs at CH$_3$.
\item \textbf{Aldehydes more reactive than ketones} — less steric hindrance and less +I reduction of C$^{\delta+}$.
\item \textbf{Carboxylic acids do not give typical C=O reactions} — resonance of --COOH reduces the electrophilicity of the C.
\item \textbf{Carboxylic acids are more acidic than phenols} — carboxylate has two equivalent resonance forms with charge on O atoms only.
\item \textbf{Chloroacetic acid stronger than acetic acid} — $-I$ of Cl stabilises the carboxylate ion.
\item \textbf{Aniline is a weaker base than ethylamine/NH$_3$} — lone pair delocalised into the ring.
\item \textbf{(CH$_3$)$_3$N weaker base than (CH$_3$)$_2$NH in water} — poorer H-bond solvation of R$_3$NH$^+$ and steric crowding.
\item \textbf{Aniline gives meta product on nitration} — in strong acid it forms anilinium ion, which is meta-directing.
\item \textbf{Aniline does not undergo Friedel--Crafts} — it forms a salt with AlCl$_3$; N$^+$ deactivates the ring.
\item \textbf{Acetylation of --NH$_2$ before nitration/bromination} — lowers activation, avoids polysubstitution and oxidation.
\item \textbf{Gabriel synthesis cannot give aryl amines} — aryl halides do not undergo nucleophilic substitution with phthalimide anion.
\item \textbf{Aryl diazonium salts are more stable than alkyl ones} — resonance with the ring.
\item \textbf{Sucrose is non-reducing} — no free anomeric --OH (C1 of glucose and C2 of fructose are bonded).
\item \textbf{Amino acids have high m.p.\ and are water-soluble} — exist as dipolar zwitterions (salt-like).
\end{fl}\par}
\end{multicols}

\clearpage
\begin{tcolorbox}[enhanced,colback=navy,colframe=navy,arc=2pt,left=6pt,right=6pt,top=2pt,bottom=2pt,before skip=4pt]
\color{white}\sffamily\bfseries\small BOARD CHECKLIST — Organic Chemistry \hfill\footnotesize\mdseries tick when practised ☐
\end{tcolorbox}
\begin{multicols}{2}\raggedcolumns
\scriptsize
\textbf{\textcolor{uVI}{Ch 6}} Classification and naming; preparation (alcohols, alkenes, Finkelstein, Swarts, Sandmeyer); S$_\text{N}$1 vs S$_\text{N}$2 with stereochemistry; chirality, racemisation; Zaitsev; Grignard, Wurtz, Wurtz--Fittig; haloarene unreactivity, o/p directing; polyhalogen uses.\\
\textbf{\textcolor{uVII}{Ch 7}} Hydration and hydroboration; phenol preparations (cumene); acidity comparisons; Lucas; dehydration mechanism; PCC vs KMnO$_4$; Kolbe, Reimer--Tiemann; bromination/nitration of phenol; Williamson (right halide); ether cleavage; anisole reactions; methanol/ethanol.\\
\textbf{\textcolor{uVIII}{Ch 8}} Rosenmund, Stephen, Etard, Gatterman--Koch, DIBAL-H; nucleophilic addition mechanism and reactivity; NH$_2$--Z derivatives; Clemmensen, Wolff--Kishner; Tollens', Fehling's, iodoform; aldol, cross aldol, Cannizzaro; acid preparation and acidity order; HVZ, decarboxylation, Kolbe electrolysis.\\
\textbf{\textcolor{uIX}{Ch 9}} Preparations (Gabriel, Hoffmann, reduction); basicity orders with reasons; carbylamine, Hinsberg, HNO$_2$; aniline substitution and protection; diazotisation, Sandmeyer/Gattermann, coupling; conversions.\\
\textbf{\textcolor{uX}{Ch 10}} Glucose open-chain evidence and its failures; pyranose/anomers; D/L; disaccharides and polysaccharides table; amino acids, zwitterion, peptide bond; protein structure levels, denaturation; enzymes; vitamins table; DNA vs RNA; hormones.
\end{multicols}

\begin{tcolorbox}[enhanced,colback=white,colframe=navy,boxrule=0.8pt,arc=2pt,title={\sffamily\bfseries\footnotesize My notes},colbacktitle=navy,coltitle=white,height fill,before skip=4pt]
\end{tcolorbox}
\end{document}
